\documentclass[11pt,a4paper]{article}

% ===================== Lingua e codifica =====================
\usepackage[utf8]{inputenc}
\usepackage[T1]{fontenc}
\usepackage[english]{babel}   % usa [italian] se scrivi in italiano

% ===================== Matematica =====================
\usepackage{amsmath,amssymb,amsthm,mathtools}
\usepackage{bm}

% ===================== Layout =====================
\usepackage[a4paper,margin=2.5cm]{geometry}
\usepackage{microtype}
\usepackage{parskip}           % opzionale: spazio tra paragrafi

% ===================== Figure e tabelle =====================
\usepackage{graphicx}
\usepackage{booktabs}
\usepackage{caption}
\usepackage{subcaption}

% ===================== Link e riferimenti =====================
\usepackage[hidelinks]{hyperref}
\usepackage{cleveref}          % va caricato DOPO hyperref e amsthm

% ===================== Bibliografia =====================
%\usepackage[backend=biber,style=alphabetic,sorting=nyt]{biblatex}
%\addbibresource{bibliografia.bib}

% ===================== Ambienti teorema =====================
\theoremstyle{plain}
\newtheorem{theorem}{Theorem}[section]
\newtheorem{lemma}[theorem]{Lemma}
\newtheorem{proposition}[theorem]{Proposition}
\newtheorem{corollary}[theorem]{Corollary}

\theoremstyle{definition}
\newtheorem{definition}[theorem]{Definition}
\newtheorem{example}[theorem]{Example}
\newtheorem{remark}[theorem]{Remark}
\newtheorem{notation}[theorem]{Notation}

% ===================== Operatori e macro =====================
\newcommand{\R}{\mathbb{R}}
\newcommand{\N}{\mathbb{N}}
\newcommand{\Z}{\mathbb{Z}}
\newcommand{\C}{\mathbb{C}}
\newcommand{\Q}{\mathbb{Q}}

\DeclareMathOperator{\Span}{span}
\DeclareMathOperator{\rank}{rank}
\DeclareMathOperator{\Ker}{ker}
\DeclareMathOperator{\Ima}{Im}
\DeclareMathOperator*{\argmin}{arg\,min}
\DeclareMathOperator*{\argmax}{arg\,max}

\newcommand{\eps}{\varepsilon}
\newcommand{\norm}[1]{\left\lVert #1 \right\rVert}
\newcommand{\abs}[1]{\left\lvert #1 \right\rvert}
\newcommand{\inner}[2]{\left\langle #1,\, #2 \right\rangle}
\newcommand{\dd}{\,\mathrm{d}}          % differenziale: \int f(x)\dd x

% ===================== Metadati =====================
\title{Un titolo chiaro e specifico\\[2pt]
       \large Un sottotitolo opzionale}
\author{Nome Cognome\thanks{Università di ..., \texttt{email@example.com}}}
\date{\today}

% ===================== Documento =====================
\begin{document}

\maketitle

\begin{abstract}
Riassunto di 100--250 parole: problema affrontato, approccio, risultato
principale, perché è significativo. Evita citazioni e formule pesanti.
\end{abstract}

\tableofcontents   % utile solo per articoli lunghi; altrimenti rimuovi

% ---------------------------------------------------------------
\section{Introduction}
\label{sec:intro}

Contesto e stato dell'arte. Motivazione del problema. Enunciato
\emph{informale} del risultato principale. Elenco dei contributi originali.

\subsection{Motivation and background}
\label{sec:motivation}

\subsection{Main results}
\label{sec:main-results-intro}

\begin{theorem}[Versione informale, cfr.\ \cref{thm:main}]
...
\end{theorem}

\subsection{Organization of the paper}
La \cref{sec:prelim} richiama le nozioni di base; la \cref{sec:main}
contiene i risultati principali; ...

% ---------------------------------------------------------------
\section{Preliminaries}
\label{sec:prelim}

Notazioni, definizioni e risultati noti (con citazioni, es.\
\cite{esempio2024}). Fissa le convenzioni una volta per tutte.

\begin{definition}\label{def:chiave}
...
\end{definition}

% ---------------------------------------------------------------
\section{Main results}
\label{sec:main}

\begin{theorem}
...
\end{theorem}

\begin{proof}
...
\end{proof}

\begin{lemma}\label{lem:aux}
...
\end{lemma}

\begin{proof}[Proof of \cref{lem:aux}]
...
\end{proof}

\begin{corollary}\label{cor:conseguenza}
Segue da \cref{thm:main} e \cref{lem:aux} che \dots
\end{corollary}

% ---------------------------------------------------------------
\section{Examples and applications}
\label{sec:examples}

Let's start by describing the typical situation.
$L$ is a first order language, $T$ is a theory for $L$ with some properties: primitive universal, irreducible, often with Q.E.
We want to study how external parameters leave a trace on a fixed model. One can prove that all results don't really depend on the fixed model, but only on $T$.
Let $M \models T$. We can add to $L$ two sorts: $V$ and $P$, standing for variable sort and parameter sort respectively. 
So, let's fix $M \models T$. 
\begin{definition}
	$S_x(M) = \{p(x) : \exists N \geq M \ \exists c \in N p(x) = tp_N(c/M)\}$.
\end{definition}

We will consider $\gamma = \{\phi_1,..,\phi_n \} \subseteq L$ a finite set of formulas in $L$ with the following two properties:
\begin{enumerate}
	\item	closed under boolean combinations
	\item	for every $i = 1,..,n$ $\phi_i$ will be of the form $\phi_i(x;y)$ where $x$ is a single variable in the sort $P$, 
		while $y$ is a vector of variables in the sort $V$   
\end{enumerate}

We denote with $L_y := \{\theta(y) : \theta(y) \in L\}$.

Ideally, $V$ will be interpreted as $M$, while $P$ will be interpreted as the disjoint union of $N \models T$ such that $N \geq M$.
So when considering $\phi(x;y) \in \gamma$, $\phi(c,V)$ is the trace that $c$ leaves on $M$ through $\phi$; more explicitly $\{m \in M : N \models \phi(c,m)\}$.
Such traces are definable using external parameters (namely $c$ in the above case) but may not be definable in $M$ using quantifier free formulas without parameters.
This depends on $c$; we are interested in those $c$ that leave a trace as close as possible to be definable.
The meaning of the last sentence is explained in what follows.

\begin{example}\label{ex:uno}
For example consider the case $L = \{<\}, T = T_{dlo}, \gamma = \{x < y, x \leq y, x > y, x \geq y\}, M = \mathbb{Q}$.
Let $N=\mathbb{R} \times 3$ and let $\forall a,b,c \in \mathbb{Q} (a,0) < (b,1) < (c,2)$.
Embed $M$ in $N$ as $\mathbb{Q} \times {1}$. Then the $\gamma$-traces left by $(0,0),(0,2)$ are definable by q.f. formulas without parameters, while those left by $(\sqrt{2},1)$ are not.
Now observe that for every $\theta(y) \in L$ and $\phi(x,y) \in \gamma$ 
\[
	N \models \forall y \ (\theta(y) \to \phi(c,y)) \iff S_{\gamma}(M) \models \mathcal{R}_{\neg \phi(x,y),\theta(y)}(tp_N(c/M))
\]

Since $T_{dlo}$ has Q.E. and every two elements have the same atomic type, we will always have $\theta(M) = M$ or $\theta(M) = \emptyset$, so we can omit $\theta$.
Consider the atomic type $\{R_{x < y}\}(z)$. It is immediate that this type is maximal since it determines every trace of every $\phi(x,y) \in \gamma$.
So there must be an element $a_{\infty} \in \mathcal{J}$ such that $\mathcal{J} \models R_{x < y}(a_{\infty})$.
$\mathcal{J}$ embeds in $S_{\gamma}(M)$ but since every trace is determined, the image of $a_{\infty}$ is determined.
In particular there cannot be two elements in $\mathcal{J}$ that model $\mathcal{R}_{x < y}$.
So we have $\mathcal{J} = \{a_{\infty}, a_{-\infty}\}$.

\end{example}

The above example is extreme, in the sense that the trace were definable or not definable without nothing between. 
In general this is not the case as the following example shows. This should illustrate even better what 'as close to be definable as possible' means.
Suppose that $\mathcal{J} \models \mathcal{R}_{\phi(x;y),\theta(y)}(a) \land \mathcal{R}_{\neg \phi(x;y), \eta(y)}(a)$ and let $f,g: \mathcal{J} \to S_{\gamma}(M)$ embeddings.
Now, $f(a)$ and $g(a)$ may be very different types, but they agree that their trace for $\phi(x,y)$ will contain $\eta(M)$ and will be disjoint from $\theta(M)$.
Since $a \in \mathcal{J}$ we know that its atomic type is maximal, which amounts to saying that $f(a)$ and $g(a)$ will differ as little as possible.

% ---------------------------------------------------------------
\section{Discussion and open problems}
\label{sec:discussion}

Limiti dell'approccio, possibili estensioni, domande aperte.

% ---------------------------------------------------------------
\section*{Acknowledgements}
Ringrazio \dots

% ---------------------------------------------------------------
%\printbibliography

% ---------------------------------------------------------------
\appendix
\section{Technical proofs}
\label{app:proofs}

Dimostrazioni lunghe o tecniche spostate qui.

\end{document}
